% -------------------------------------------------------
\section{Сквозной анализ: три парадигмы, девять моделей}
% -------------------------------------------------------

\subsection{Сравнительная таблица результатов по всем экспериментам}

Таблица~\ref{tab:cross_experiment} представляет агрегированные результаты по всем трём экспериментам для каждой из девяти моделей.

\begin{table}[h!]
\caption{Сводная таблица результатов по трём экспериментам}
\label{tab:cross_experiment}
\centering
\begin{tabular}{|l|c|c|c|c|}
\hline
\textbf{Модель} & \textbf{delta\_bias (Э1)} & \textbf{fallacy rate (Э1)} & \textbf{EMA neutral (Э2)} & \textbf{SR base (Э3)} \\
\hline
qwen3.5-397b-a17b & $-0{,}005$ & $0{,}2\%$  & 0,764 & 0,22 \\
glm-4.7-flash     & $+0{,}028$ & $0{,}8\%$  & 0,332 & 0,04 \\
gemma-3-12b-it    & $+0{,}039$ & $2{,}9\%$  & 0,028 & 0,06 \\
claude-sonnet-4.6 & $+0{,}105$ & $0{,}2\%$  & 0,790 & 0,28 \\
deepseek-v3.2     & $+0{,}200$ & $18{,}7\%$ & 0,658 & 0,10 \\
grok-4.1-fast     & $+0{,}249$ & $22{,}9\%$ & 0,424 & 0,14 \\
gpt-5.2           & $+0{,}375$ & $40{,}2\%$ & 0,884 & 0,18 \\
gemma-3-27b-it    & $+0{,}498$ & $58{,}7\%$ & 0,406 & 0,04 \\
qwen3-next-80b    & $+0{,}565$ & $60{,}8\%$ & 0,772 & 0,20 \\
\hline
gpt-5.2 reasoning & ---        & ---        & ---   & 0,38 \\
\hline
\end{tabular}
\end{table}

\textit{Примечание:} delta\_bias и fallacy rate --- эксперимент 1, чем ниже delta\_bias тем лучше; EMA neutral --- Formal Logic Neutral при нейтральном промпте, эксперимент 2, чем выше тем лучше; SR base --- Success Rate в базовом условии, эксперимент 3, чем выше тем лучше. Таблица отсортирована по delta\_bias. \texttt{gpt-5.2} reasoning выделен отдельно как режим с принципиально иным типом вывода.

\subsection{Проверка гипотез о масштабировании и внутрисемейный парадокс}

Гипотеза H2 о снижении предубеждений с ростом модели получила частичное подтверждение. На уровне общей тенденции зависимость наблюдается: более мощные модели (\texttt{gpt-5.2}, \texttt{claude-sonnet-4.6}) систематически демонстрируют более низкий уровень предубеждений во всех трёх парадигмах. Однако внутри семейств картина неоднородна.

В семействе Gemma направление эффекта противоречит ожидаемому: \texttt{gemma-3-27b-it} демонстрирует несравнимо более высокий токен-bias (delta\_bias $= 0{,}498$, fallacy rate $58{,}7\%$), чем \texttt{gemma-3-12b-it} (delta\_bias $= 0{,}039$, fallacy rate $2{,}9\%$). Это является наиболее ярким проявлением внутрисемейного парадокса масштаба, зафиксированного в эксперименте 1. В семействе Qwen наблюдается противоположный паттерн: \texttt{qwen3-next-80b} при меньшем числе активных параметров демонстрирует существенно более высокий bias (delta\_bias $= 0{,}565$), чем \texttt{qwen3.5-397b-a17b} (delta\_bias $= -0{,}005$). Оба семейства указывают на то, что масштаб сам по себе не является надёжным предиктором устойчивости к когнитивным предубеждениям, и подчёркивают роль post-training alignment в формировании профиля предубеждений~\cite{batista2026rational}.

\textbf{Формальная проверка.} Для количественной оценки связи между масштабом и предубеждениями проведён ранговый корреляционный анализ Спирмена по шести открытым моделям с верифицируемым числом активных параметров: \texttt{gemma-3-12b-it} (12B), \texttt{qwen3.5-397b-a17b} (17B active), \texttt{gemma-3-27b-it} (27B), \texttt{glm-4.7-flash} (30B), \texttt{deepseek-v3.2} (37B active), \texttt{qwen3-next-80b} (80B). Проприетарные модели (\texttt{gpt-5.2}, \texttt{claude-sonnet-4.6}, \texttt{grok-4.1-fast}) исключены из анализа из-за отсутствия верифицируемых архитектурных данных; IKP-оценки~\cite{li2026ikp} дают лишь нижние границы с медианной погрешностью $1{,}59\times$, недостаточные для корреляционного анализа.

Результаты: $\rho_S$(active params, delta\_bias) $= 0{,}600$, $p = 0{,}208$; $\rho_S$(active params, EMA) $= 0{,}657$, $p = 0{,}156$; $\rho_S$(active params, SR) $= 0{,}116$, $p = 0{,}827$. Ни одна из корреляций не достигает значимости при $\alpha = 0{,}05$. Это формально подтверждает качественный вывод: при контроле по открытым моделям число активных параметров не является статистически значимым предиктором ни вероятностного предубеждения, ни дедуктивной компетентности, ни индуктивного успеха. Единственная умеренная положительная тенденция наблюдается для EMA ($\rho = 0{,}657$), однако она полностью обусловлена разницей между \texttt{gemma-3-12b-it} (EMA $= 0{,}106$) и остальными моделями и исчезает при исключении этой точки.

\subsection{Проверка гипотезы о спектральной природе предубеждений}

Гипотеза H3 подтверждена: уровень предубеждений систематически возрастает от эксперимента 1 к эксперименту 3. В задаче Линды в uncorrelated-условии ошибок нет ни у одной модели (все $n_{12} = 0$); в задаче Уэйсона EMA в нейтральном условии при нейтральном промпте не превышает 0,884 даже у лучшей модели, а у \texttt{gemma-3-12b-it} падает до 0,028; в задаче 2-4-6 CTR $> 0{,}60$ у всех моделей без исключения. Таким образом, индуктивное открытие правил остаётся наиболее уязвимой к предубеждению подтверждения областью для современных LLM, что согласуется с выводами Banatt и соавторов~\cite{banatt2024wilt} о тотальном доминировании стратегии позитивного тестирования, и с данными CogBench~\cite{coda2024cogbench} о систематическом разрыве между LLM и нормативными агентами именно в задачах фальсифицирующего поиска.

\subsection{Итоги: Proof-of-Concept фреймворка оценки когнитивных предубеждений}

Совокупность результатов трёх экспериментов позволяет сформулировать следующие итоговые положения.

Во-первых, когнитивные предубеждения в LLM являются устойчивым, воспроизводимым феноменом. Токен-bias зафиксирован у восьми из девяти моделей в эксперименте 1; предубеждение подтверждения в дедуктивном рассуждении --- у всех девяти в эксперименте 2; CTR $\geq 0{,}629$ у всех non-reasoning моделей в эксперименте 3. Гипотеза H1 подтверждена с оговоркой: \texttt{qwen3.5-397b-a17b} не проявляет значимого токен-bias в задаче Линды, однако демонстрирует выраженное предубеждение подтверждения в задачах Уэйсона и 2-4-6 (CTR $= 0{,}850$).

Во-вторых, профили предубеждений по экспериментам не коррелируют монотонно. Наглядный пример: \texttt{gpt-5.2} имеет высокий токен-bias в задаче Линды (fallacy rate $40{,}2\%$, delta\_bias $= 0{,}375$), но лучшую EMA в задаче Уэйсона ($0{,}884$) и наилучший SR в задаче 2-4-6 среди non-reasoning моделей ($0{,}18$, а в reasoning-режиме --- $0{,}38$). Это означает, что уязвимость к разным типам предубеждений не является единым латентным фактором, и ни одна агрегированная метрика не заменяет полного профиля.

В-третьих, масштаб модели не является надёжным предиктором устойчивости к предубеждениям, что согласуется с эффектами обратного масштабирования, описанными Dasgupta и соавторами~\cite{dasgupta2024language}. Внутрисемейный парадокс Gemma проявляется по-разному в разных экспериментах: в задаче Линды \texttt{gemma-3-27b-it} значительно хуже, чем \texttt{gemma-3-12b-it}; в задаче Уэйсона --- значительно лучше (EMA $0{,}578$ vs $0{,}106$); в задаче 2-4-6 обе модели слабы (SR $0{,}04$ и $0{,}06$). Это подчёркивает, что post-training alignment оказывает неоднородное воздействие на разные типы рассуждения~\cite{batista2026rational}.

В-четвёртых, reasoning является единственным вмешательством, дающим качественный скачок: SR \texttt{gpt-5.2} reasoning + adaptive ($0{,}48$) вдвое превышает лучший non-reasoning результат ($0{,}28$). Prompt engineering в форме CoT улучшает результаты только у моделей с достаточной базовой компетентностью и может давать обратный эффект у слабых моделей~\cite{kojima2022large}.
